\section*{अध्याय \ref{ch:g6:measurement-in-science} --- विज्ञान में मापन: मात्रक और उपकरण}

\begin{solution}{exo:g6:measurement-in-science:1}
$2$ \unit{m}; \ $8$ \unit{g}; \ $55$ \unit{min}; \ $39$
\unit{\celsius}।
\end{solution}

\begin{solution}{exo:g6:measurement-in-science:2}
\qty{320}{cm}; \ \qty{0.75}{kg}; \ \qty{240}{s}; \ \qty{5.6}{km}।
\end{solution}

\begin{solution}{exo:g6:measurement-in-science:3}
उसके पायदान $10$ के नहीं, $60$ के हैं, इसलिए दशमलव की सीढ़ी झूठ बोल
देती है: $\qty{3.5}{min} = 3 \times 60 + 30 = \qty{210}{s}$, जबकि
सीढ़ी ग़लती से \qty{350}{s} देती।
\end{solution}

\begin{solution}{exo:g6:measurement-in-science:4}
हर पायदान: $200 \div 5 = \qty{40}{mL}$। $300$ से दो पायदान ऊपर:
$300 + 80 = \qty{380}{mL}$।
\end{solution}

\begin{solution}{exo:g6:measurement-in-science:5}
ख़मीर: रसोई की तराज़ू (\omterm{def:g3:measuring-mass:units}{ग्राम} वाले पायदान --- वज़न वाली तराज़ू चुटकी भर
को पकड़ ही नहीं पाती)। नाड़ी: स्टॉपवॉच (सेकंड और उससे बारीक पायदान)।
पन्ने की चौड़ाई: रूलर (मिलीमीटर वाले पायदान)। स्कूल का मैदान: लंबा
मापक फ़ीता या सर्वेक्षक का पहिया (मीटरों वाला परास)।
\end{solution}

\begin{solution}{exo:g6:measurement-in-science:6}
परखने वाला क़दम: पेंसिल का अंदाज़ा \qty{20}{cm} से कम का होता है,
इसलिए \qty{17.4}{m} पहली ही नज़र में समझदारी की जाँच में फ़ेल हो जाता
है। उसने \omterm{def:g2:measuring-length:units}{सेंटीमीटर} पढ़े और \omterm{def:g2:measuring-length:units}{मीटर} लिख दिए: असली मान \qty{17.4}{cm} है।
\end{solution}

\begin{solution}{exo:g6:measurement-in-science:7}
ईमानदार: ``$18$ और \qty{19}{\celsius} के बीच, $19$ के ज़्यादा
पास'' (या ``लगभग \qty{18.8}{\celsius}'')। बेईमान:
``\qty{18.7643}{\celsius}'' --- यानी वह परिशुद्धि जिसका वादा निशान
कर ही नहीं सकते।
\end{solution}

\begin{solution}{exo:g6:measurement-in-science:8}
किसी की ग़लती नहीं है: उँगलियाँ बाल भर के फ़र्क़ से चालू और बंद होती
हैं, और थोड़ा बिखराव मापन की सामान्य साँस है। समझदारी वाली रिपोर्ट:
पाठ \qty{6.2}{s} के आस-पास जुटते हैं --- ``लगभग \qty{6.2}{s}''।
\end{solution}

\begin{solution}{exo:g6:measurement-in-science:9}
\qty{450}{g} का चौथाई: लगभग \qty{112}{g} (\qty{113}{g} भी ठीक है)।
यह काम दिखाता है कि साझा \omterm{def:g6:measurement-in-science:measuring}{मात्रक} क्यों ज़रूरी हैं: एक देश के मात्रकों
में लिखी पाक-विधि को पहले बदलना पड़ता है, तभी दूसरे देश के उपकरण उसका
कहा मान सकते हैं।
\end{solution}

\begin{solution}{exo:g6:measurement-in-science:10}
टिकट $1$ के निशान से $3.4$ के निशान तक फैला है: $3.4 - 1 =
\qty{2.4}{cm}$। आदत यह है: किसी उपकरण के किनारे पर कभी भरोसा मत करो
--- कोई चुना हुआ निशान लेकर शुरू करो और घटा लो, फिर घिसे हुए रूलर
अपनी ताक़त खो देते हैं।
\end{solution}

\begin{solution}{exo:g6:measurement-in-science:11}
दोनों पाठ \qty{0.10}{m} से बिखरे हैं, इसलिए किसी भी एक का आख़िरी अंक
भरोसे लायक नहीं; वे $18.3$ के आस-पास जुटते हैं।
``\qty{18.3}{m}'' लिखने से हर वह अंक बचा रहता है जिसका वादा फ़ीता
ईमानदारी से कर सकता है, और उससे एक भी ज़्यादा नहीं।
\end{solution}

\begin{solution}{exo:g6:measurement-in-science:12}
आफ़तें: किन्हीं दो क़स्बों के ``पैर'' आपस में नहीं मिलते, इसलिए एक
क़स्बे में कटे शहतीर दूसरे क़स्बे के घरों में कभी नहीं बैठते; और बेटे
के गद्दी पर आते ही ख़ुद \omterm{def:g6:measurement-in-science:measuring}{मात्रक} की लंबाई बदल जाती है --- राजा के साथ
राज्य का हर मापन ख़त्म हो जाता है। धातु की छड़ ने दोनों ठीक किए: एक न
बदलने वाली लंबाई, जिसकी नक़लें हर क़स्बे को दी जा सकें और जो हर राजा
से ज़्यादा जिए। उसकी अपनी कमज़ोरी: कोई अकेली वस्तु ख़रोंची जा सकती है,
चुराई जा सकती है, या धीरे-धीरे बदल सकती है, और नक़लें जाँचने के लिए
पूरी दुनिया को उसी तहख़ाने तक जाना पड़ता था --- जब तक \omterm{def:g6:measurement-in-science:measuring}{मात्रक} प्रकृति
के नियतांकों पर फिर से नहीं गढ़े गए, जो न घिसते हैं और एक साथ हर जगह
ड्यूटी पर रहते हैं।
\end{solution}

\begin{solution}{pb:g6:measurement-in-science:1}
\textbf{1.} हर तौल उस पाठ में बोझ जोड़ देती है जो सुई पहले से दिखा
रही थी; शून्य से हटी शुरुआत पूरे दिन के हर पाठ को ज़हरीला कर देती है
--- इसलिए निष्पक्ष शुरुआत की जाँच सबसे पहले होनी चाहिए।
\textbf{2.} $500 + 200 = \qty{700}{g}$।
\textbf{3.} वह \qty{20}{g} कम बताती है। सुई के \qty{700}{g} कहने
तक आटा डालते जाने पर नानबाई असल में \qty{720}{g} डाल चुकी होती है:
तराज़ू नानबाई से ग्राहक के दाम से ज़्यादा सामान दिलवाती है --- यानी वह
ग्राहकों की चापलूसी करती है, और हर सौदे में नानबाई गँवाती है।
\textbf{4.} $500 + 200 + 50 = \qty{750}{g}$।
\textbf{5.} नहीं: \qty{2}{g} का टिकट तराज़ू के \qty{5}{g} वाले
पायदानों के नीचे ग़ायब हो जाता है --- वह उसके इलाक़े से बाहर है।
\textbf{6.} हर पायदान $200 \div 4 = \qty{50}{mL}$ का।
\textbf{7.} पानी $600 + 50 = \qty{650}{mL}$ पर है; ``लगभग
\qty{700}{mL}'' पूरा एक पायदान गोल कर देता है --- इतने अच्छे उपकरण
के लिए बहुत ढीली रिपोर्ट।
\textbf{8.} आँख निशान के बराबर हो: नीचे से पढ़ने पर पानी की सतह
खिसकी हुई दिखती है --- ``ठीक सामने से'' वाला नियम टूटा।
\textbf{9.} \qty{10}{mL} का भटकाव --- यानी उसके अपने \qty{50}{mL}
वाले पायदान से बेहतर (छोटा): स्कूल के जग के लिए ठीक है।
\textbf{10.} तेज़, $2$ मिनट से।
\textbf{11.} हफ़्ते भर बाद $2 \times 7 = 14$ मिनट आगे।
\textbf{12.} पहली घड़ी: $2$ मिनट पीछे करो; दूसरी: $2$ मिनट आगे;
तीसरी: कुछ नहीं --- वह सही है। (तीनों का औसत निकालने से कुछ ठीक नहीं
होता और अकेली ईमानदार घड़ी झूठी पड़ जाती है।)
\textbf{13.} जैसे: \emph{शून्य}: किसी उपकरण की किसी भी बात पर भरोसा
करने से पहले जाँच लो कि ख़ाली हालत में वह शून्य दिखाता है।
\emph{पढ़ना}: पायदान का मान सीखो और ठीक सामने से पढ़ो, और निशानों के
बीच ईमानदार विनम्रता रखो। \emph{\omterm{def:g6:measurement-in-science:measuring}{मात्रक}}: \omterm{def:g6:measurement-in-science:measuring}{मात्रक} के बिना संख्या मापन
है ही नहीं --- \omterm{def:g6:measurement-in-science:measuring}{मात्रक} हमेशा लिखो।
\end{solution}
